% !TeX root = ../main.tex
\section*{Introducción}
La muerte del testigo fiel constituye una de las paradojas decisivas del Apocalipsis: aquello que el poder de la bestia presenta como derrota, la visión de Juan lo contempla desde la victoria del Cordero.
Sin embargo, cuando el profeta anuncia que algunos muertos \enquote{vivieron y reinaron con Cristo} y llama a esto \enquote{la primera resurrección} (Ap 20:4--6), surge un problema hermenéutico que ha dividido a la iglesia: ¿se refiere a la restauración corporal de los santos antes de la consumación o al ingreso de los creyentes fallecidos en la vida celestial?
La respuesta no puede depender simplemente de que el texto sea simbólico, pues las visiones apocalípticas comunican realidades precisas mediante sus imágenes.

Las confesiones reformadas no definen expresamente la \enquote{primera resurrección}; sí afirman que las almas de los justos están con Cristo después de morir y esperan la resurrección corporal general del último día.
Ello establece un horizonte doctrinal, pero no sustituye el examen del pasaje.
Este ensayo sostendrá que, en Apocalipsis 20:4--6, la \enquote{primera resurrección} no designa una resurrección corporal anterior a la consumación ni, primariamente, la regeneración del creyente, sino la entrada victoriosa de los fieles muertos, representados paradigmáticamente por los mártires, en la vida y el reinado celestial con Cristo durante la era interadvental; la resurrección corporal de justos e injustos permanece como un acontecimiento general del último día.
La tesis se defenderá atendiendo a la recapitulación narrativa, al contraste entre las lecturas amileniales, a las correspondencias internas del Apocalipsis y a las objeciones de la interpretación premilenial.
